\section{结论}
本章中，我们探索了使用非平稳无记忆策略解决状态观测或动作执行中常值延迟的可能性。
确实，无记忆策略由于对增广状态的部分观测，自然会导致动态的非平稳性。
值得注意的是，在非平稳和马尔可夫策略空间中存在最优无记忆策略 \cite[Theorem~5.1]{derman2021acting}。

基于这一结果，我们设计了一种能够适应过程回合内非平稳性的方法。
受考虑此类场景的终身学习文献启发，我们提出学习一个超策略，它以时间为输入，采样在该时间被查询的策略的参数。
为了优化未来性能，我们为其设计了一个估计量，其基础依赖于平滑非平稳性的假设。
具体来说，该估计量重用过去样本来估计未来，并随着样本变旧对其进行指数折扣。
在平滑条件下，我们证明了估计量的偏差是有界的，并且指数折扣可以在一定程度上控制它。
值得注意的是，当环境和超策略都是平稳的时，偏差消失。
在未来性能估计量之上，目标函数增加了两项。
第一，添加了新超参数集下过去性能的估计。
这一项通过迫使超策略在过去样本上表现良好来对抗灾难性遗忘。
第二，增加了对方差的惩罚。
它禁止超策略的过度非平稳性，这种非平稳性将是过拟合环境过去非平稳性和未来泛化不良的症状。
通过梯度优化来优化这一目标得到了我们的算法 POLIS。

利用该算法，我们提出了一个简单修改以考虑延迟。
与 dSARSA 类似 \cite{schuitema2010control}，我们在性能估计期间将奖励在时间上向后偏移，以重新对齐所选策略与其结果。

随后，我们展示了 POLIS 在不同终身场景（有延迟和无延迟）中的实证评估。
该算法展示了其能够有效学习非平稳性的结构以适应未来。
POLIS 还避免了不需要的额外非平稳性，这主要得益于目标函数中的惩罚项。
在延迟任务上测试时，POLIS 随着延迟增加表现出鲁棒的行为。
然而，与水坝环境中的平稳超策略相比，POLIS 的回报有时具有较高的方差。

未来研究可以考虑随机延迟。
在这种情况下，无记忆策略的最大优势在于其输入是 $\statespace$ 中的一个状态，并且不像增广方法那样依赖于延迟。
因此，无记忆策略具有小而固定的输入维度 $\cardinal{\statespace}$。
